
\chapter{\label{III_1} Une intelligence hybride}

L'intégration de l'apprentissage automatique aux processus d'archivage audiovisuel,relève d'un large potentiel de découvrabilité du contenu de ces archives. Toutefois, ces outils restent difficiles à mettre en application au sein d'un tel milieu, nécessitant la mise en place d'efforts communs afin de "redéfinir en permanence la répartition du travail en l'homme et la machine dans le sens d'un partenariat" \footcite{zotero-item-584}. En effet, le processus tel que nous l'avons décrit ne saurait être entièrement automatisé, et laisserait plutôt la place à une cohabitation entre l'extraction automatique et l'indexation documentaire manuelle.

Cette cohabitation est illustrée par le concept d'"\gls{Intelhyb}" définie par Zeynep Akata comme :

\begin{quote}
	La combinaison de l'intelligence humaine et de l'intelligence des machines qui complète l'intellect et les capacités humaines, plutôt que de les remplacer, afin de prendre des décisions judicieuses, d'effectuer des actions appropriées et d'atteindre des objectifs qui ne peuvent être atteints ni par les humains ni par les machines seules [Traduction personnelle]" \footnote{\textit{"the combination of human and machine intelligence, augmenting human intellect and capabilities instead of replacing them and achieving goals that were unreachable by either humans or machines."} \cite{akata2020}}.
\end{quote}
Ce concept implique notamment la mise en place de quatres étapes distinctes : la collaboration entre l'homme et la machine, l'adaptabilité, la responsabilité, et l'explicabilité \footcite{akata2020}.



\section{\label{III_1_A}Une segmentation manuelle} 

La préparation des pratiques utilisateurs à ces nouvelles fonctionnalités passe notamment par une préparation méthodologique de l'interface, dédiée aux collections audiovisuelles. Par exemple, la mise en place d'un outil de visionnage et d'annotation vidéo pourrait permettre de préparer les utilisateurs à de telles pratiques. Cet outil pourrait répondre à un double enjeu, celui d'encourager les pratiques d'annotation et de segmentation sur les fichers numériques, ainsi que celui d'éventuellement collecter les données produites par un tel outil afin d'entraîner les modèles à implémenter. (collection as data)
La mise en place de cet outil permettrait également de faciliter le contrôle qualité, en comparant les données manuelles et automatiques. 

L'outil Mediascope, développé par l'\gls{ina} à destination des chercheurs à l'INAthèque, est un outil d'aide à l'analyse des contenus audiovisuels. Il permet de structurer le développement temporel du document à l'aide d'annotations associées à des "time-codes" et à des images fixes \footcite{zotero-item-746}. L'utilisateur peut ainsi visionner un document audiovisuel, annoter des marqueurs temporels sur des images fixes, ainsi qu'y saisir du texte. L'outil permet également l'export des images, du texte, et de statistiques sur le contenu \footcite{zotero-item-746}. Mediascope unifie ainsi la lecture vidéo et l'indexation d'annotations au sein d'une même interface.

Cette proposition s'inscrit dans un document global d'améliorations : cf annexe :

Au sein de l'application Skinsoft, la gestion des médias se limite à la visualisation simple d'un fichier mutimédia lié à une notice. Afin d'améliorer cette visualisation des médias, nous avons proposé au cours de notre stage, sur l'insipration du Médiascope, une visionneuse permettant le séquençage manuel et l'annotation de données. 
Cette amélioration s'appuierait d'une part sur deux possibilités de séquençage manuel. L’utilisateur peut saisir des repères temporels sur un fichier vidéo sans pour autant le modifier. Ces repères sont reversés dans des champs prévus pour accueillir les “time codes” dans la notice de l’objet numérique. Au sein du fichier vidéo, les marqueurs sont matérialisés et visibles sur la barre de défilement de la vidéo, permettant une navigation interactive. De plus, la visionneuse permettrait à l’utilisateur de segmenter le fichier vidéo en plusieurs séquences distinctes, à partir de repères temporels placés (développement en annexe). Chaque séquence identifiée est reversée dans un fichier d'item numérique propre, lié à une notice descriptive du niveau Fragment. 
D'autre part, cette amélioration permettrait deux niveaux d'annotations sur le média. Une annotation globale au niveau de la séquence qui prendrait par exemple la forme d'un résumé, de mots-clés généraux thématiques, inscrits pau sein de marqueurs temporels. Une annotation précise de tags apposés sur des éléments de l'image tels que des objets, personnes, éléments de décor, similaire à la détection d'objet prévue par \gls{yolo}. Ces éléments seraient indexés dans un champ dédié avec un ancrage temporel et spatial. (développement en annexe)
À la manière du médiascope, cette amélioration du visionnement des médias permettrait de combiner visualisation, segmentation et annotation au sein d'un même outil. 

La mise en place d'un tel outil permettrait déjà d'améliorer la recherche en la faisant porter sur des éléments précis portés par le contenu des images. De plus, le séquençage et l'annotation manuelle sur les documents audiovisuels prépare es utilisateurs à la structure des données engendrées en masse par les algorithmes à implémenter. Enfin, la production manuelle de ces données (time codes et mots-clés) faciliterait l'entraînement et le réentraînement des modèles, ainsi que le contrôle qualité effectué par l'utilisateur. Les données produites en amont permettraient en effet de constituer un jeu de données d'entraînement pour les modèles implémentés. Ces données serviraient également de référence afin d'inspecter la qualité de celles produites automatiquement par les modèles.
\section{\label{III_1_B}Visualisation du FRBR en graphe pour préparer la mise en place d’une nouvelle entité : la séquence et de nouvelles données produites automatiquement : alourdit les bdd donc nécessite des visualisations}

Visualisation des données et contrôle qualité (rappeller la sous partie des problèmes de navigation au sein du FRBR)

Comme nous l'avons évoqué dans la première partie, bien que le modèle FRBR soit élaboré afin d'optimiser la découvrabilité des collections, il complexifie toutefois la mise en place d'une interface intuitive et navigable, enjeu au coeur des préoccupations d'un éditeur tel que Skinsoft. L'implémentation de nouvelles fonctionnalités automatisées accentue cette complexité, avec la mise en place d'une nouvelle entité, celle du Fragment d'oeuvre, et la génération massive de données supplémentaires, qui risque de rendre moins claires les relations entre les notices. Pour ne pas entraver la lisibilité des liens entre les niveaux de description, il est nécessaire de proposer une vue d'ensemble de la structure des données au sein de S-Movies. 
Au cours de notre stage, nous avons proposé une visualisation de la structure de données sous forme de graphe relationnel. Cette proposition s'inspire notamment des travaux de la Bibliothèque Nationale Allemande qui ont amener à développer une visualisation en graphe du modèle FRBR afin de structurer les résultats de recherche indexés \footcite{zotero-item-749}. En s'appuyant sur la plateforme \textit{Culture Graph} et des outils de spatialisation tels que \textit{Gephi}, des algorithmes regroupent les œuvres sous des clés uniques et cartographient leurs liens sémantiques \footcite{zotero-item-749}. De plus, des mots-clés reliés aux entités sont également modélisés sous forme d'arrêtes \footcite{zotero-item-749}. Appliqués à une structure de données audiovisuelles, le même type de visualisation en graphe peut être imaginée, traduisant visuellement la densité des relations entre les entités du modèle. À terme, après l'implémentation des fonctionnalités automatiques d'indexation, les mots-clés générés et les times codes pourront, de la même manière, alimenter le graphe afin d'améliorer le contrôle qualité sur les données. La visualisation en graphes représente ainsi un "complément précieux" aux résultats de recherche traditionnels, permettant à l'utilisateur de comprendre beaucoup plus facilement le regroupement d'entités \footcite{zotero-item-749}. La visualisation en graphe n'est pas seulement utile pour l'amélioration de la visualisation des liens entre les entités dans les résultats de recherche, mais elle permet également d'évaluer la cohérence du modèle de données dans son ensemble. En effet, "l’impression visuelle donnée par le graphe d’une œuvre est très utile car elle attire l’ attention sur les sommets qui sont cruciaux pour la structure d’un cluster d’œuvres."\footcite{zotero-item-749}. identifier d'un coup d'œil les erreurs d'association ou les ruptures de cohérence générées par les traitements automatiques avant leur validation définitive. 
Concrètement, au sein de S-Movies, nous proposons l'intégration d'une vue d'ensemble du modèle en graphe, jouxtant la vue des résultats de recherche classique. Cette vue pourra être étendue au sien des ntoices individuelles, permettant d'obtenir une granularité différente, à partir d'un item individuel. Cette vue en graphe permettrait de faciliter la navigation entre les entités descriptives et de vérifier la cohérence des liens, dans le cadre de l'implémentation d'une indexation automatisée. 
Une description détaillée de cette proposition est présentée en annexe.

Intéropérabilité des données ?
\section{\label{III_1_C} Explicabilité et transparence}  

Le contrôle de la structure et de la qualité des données dépend également de la responsabilité de l'éditeur logiciel via-à-vis d'une transparence et d'une explicabilité du processus mis en place. Comme le rapelle Jonas Arnold, la gestion de la qualité de l'indexation dépend des utilisateurs, ainsi "la transparence envers ces derniers est primordiale." \footcite{zotero-item-584}. Afin de garantir cette transparence, les processus automatiques implémentés pourront être expliqués à deux niveaux au sein de l'interface. 

Tout d'abord, via une note d'information en amont du déclenchement du processus, au sein du diagramme de décision schématisé précédemment (\ref{fig:validation}). Avant le déclenchement du processus, l'interface devra indiquer que la fonctionnalité utilisée est fondée sur l'intelligence artificielle. Le même principe est prévu pour le développement de la recherche en langage naturel, via une infobulle qui en explique brièvement le fonctionnement à des fins de transparence (Voir Annexe B, \cref{figinfobulle}). De plus, une documentation pourra être mise à la disposition des utilisateurs sous la forme d'un mode d'emploi explicatif de la fonctionnalité. Skinsoft prévoit déjà ce type de documentation à l'usage des utilisateurs métier, via un manuel \textit{wiki} dédié à chaque espace de travail ets es fonctionnalités. Le mode d'emploi pourra mentionner de manière synthétique les critères d'extraction du modèle utilisé, justifier le choix des outils ainsi que la mention de leurs limites, pouvant guider les utilisateurs dans leur contrôle qualité des données. Il serait pertinent de mentionner également les limites liées à l'entraînement du modèle, qui apprend à extraire des informations brutes, mais qu'il n'apprend pas à les placer au sein d'une structure de données, en lien avec des normes métier. 

Il s'agira également de mentionner le périmètre au sein duquel sont analysées les données, qui est celui de l'éditeur seul. En effet, l'une des réticences majeures identifiés face à l'implémentation de fonctionnalités automatiques réside dans la perte de souveraineté des données patrimoniales lors de leur injection au sein de solutions commerciales externes \footcite{zotero-item-584}. Pour un éditeur comme Skinsoft, l'enjeu consiste à garantir une architecture logicielle éthique, étanche et souveraine, au sein de laquelle les données du client ne sont jamais réutilisées à son insu pour enrichir des réservoirs tiers.


L'"\gls{Intelhyb}" telle que définie par Akata se concrétise ainsi à travers la mise en place de trois étapes intégrées à la validation du processus automatique : l'indexation du nouveau statut de données "fluides" à travers des métadonnées de processus, le contrôle de la qualité de ces données au sein de la structure logicielle, et l'explicabilité du processus afin d'assurer sa transparence \footcites{akata2020}{zotero-item-584}. Le développement actuel de \gls{llm} par Skinsoft pourra être mis à profit afin d'assister l'utilisateur dans la mise en place d'un contrôle critique. Toutefois, la mise en place de cette "\gls{Intelhyb}" au profit de la qualité des données exige une innovation de l'environnement de travail lui-même \footcite{akata2020}. Il convient ainsi d'étudier les évolutions requises au sein de S-Movies, afin d'assurer la fluidité du contrôle des données automatisées. 